\section{Discussion}\label{sec:disc}
\paragraph{What holds up.} Across all three studies, the robust result is the open-world \emph{posterior}. Keeping $h_0$ lets the planner say ``none of the above'', and it says so far more reliably than a threshold on closed-world uncertainty. In simulation this held in every scenario. On real data it held across three noise seeds and three looks at the test set, retrospective and causal. Closed-world planners are confidently wrong on most genuinely unmodeled transients; the open-world planner on a small fraction.

\paragraph{Two ways to say ``unknown''.} The Bayesian decision and the conformal one fail differently. The Bayesian $h_0$ made the fewest confident errors on real unknowns but has no false-alarm guarantee. The conformal threshold certified its false-alarm rate on fresh supernovae and, with lookahead, raised supernova accuracy above the closed-world planner's, at the price of more confident errors on unknowns. It also works on top of any planner, including a closed-world one. The choice between them is a choice of which error is costlier, and both beat closed-world planning on the unknowns.

\paragraph{What does not.} Two parts of the proposed utility did not earn their place. The adequacy term $\beta I(Z;Y_a)$, which targets measurements that test whether $h_0$ is true, made no measurable difference. An earlier toy study with retrospective sampling already suggested that planning with $h_0$ adds little beyond having it in the posterior. The decay term helped, but mainly by stopping greedy, myopic information gain from jumping to a late, high-contrast epoch and forfeiting the rest of its budget. Its benefit to fast classes was carried by kilonovae, reversed for flares (whose signal is gone before any epoch at which the proxy can act), and reversed when the forward models were wrong. And the planner was not the fastest identifier: a cheap disagreement heuristic identified the events it could identify sooner. A one-step utility has no notion of a horizon. The simple lookahead term of Section~\ref{sec:conf} used more of the budget and raised accuracy, but did not reliably improve the detection of unknowns; principled non-myopic planning remains open.

\paragraph{What the time-ordering error taught us.} Letting planners measure backwards in time did not bias the comparison between strategies, since every strategy had the same freedom, but it changed the conclusions. Retrospectively, all planners used their full budget; causally, adaptive planners used about half. With fewer measurements every posterior is less confident. That removed confident errors on both known and unknown events, turned the open-world planner's failed supernova-cost hypothesis (H3) into a pass, and cut the share of unknowns it flagged from \RealRFlagP\ to \RealCFlagP. A planner that can revisit the past cannot be evaluated as a follow-up planner. We recommend that benchmarks for sequential follow-up check the time order of every chosen sequence explicitly.

\paragraph{Limitations.} The simulated classes are analytic shapes chosen to match qualitative timescales, not published models, and amplitude is their only physical parameter, so H2 measures information about amplitude only. The decay term is a squared-signal-to-noise proxy, not the per-class information decay it stands for. There is no slew or exposure cost, a single band in simulation, geometric visibility only (no link budget), and a centralized planner, so the distributed-allocation hypothesis of the original proposal is untested. The decay weight was selected at the edge of its validation grid. On real data, measurements are resampled from interpolated light curves rather than new photons; the noise model is assumed; simultaneous two-band imaging is assumed; and the test events had been examined twice before the causal run, which we declared in advance. The conformal study's unknown-class claims reuse those test events once more, and its guarantee assumes calibration and new supernovae are exchangeable, which holds by construction for the sealed set but only approximately for the test split.

\section{Conclusion}
For follow-up of fading transients, keep an explicit ``none of the above'' hypothesis in the posterior: it is the part of the method that consistently prevents confident misclassification of the unknown, in simulation and on real ZTF events, at little or no cost on ordinary events once measurements are causal. Where a guaranteed false-alarm rate matters more than the fewest confident errors on unknowns, a conformal threshold on the same likelihoods provides it. Do not expect a one-step utility to manage urgency; the decay term we tested helps only partly and unevenly, and the planner's identification speed did not beat a simple heuristic. Finally, check the time arrow.
